\documentclass{article}
\usepackage[T1]{fontenc}
\usepackage{lmodern}
\usepackage{graphicx}
\usepackage{amsmath,amssymb}
\usepackage[a4paper,margin=2cm]{geometry}
\usepackage[absolute,overlay]{textpos}
\setlength{\TPHorizModule}{1mm}
\setlength{\TPVertModule}{1mm}
\usepackage[indonesian]{babel}
\usepackage{hyperref}
\setlength{\emergencystretch}{2em}
% unit: Halaman 90

\begin{document}

% === Halaman 90 (python/native) ===

ChaCha20

• Chacha20 merupakan modifikasi dari Salsa20 yang dipublikasikan oleh Bernstein 

pada tahun 2008.  ChaCha20 merupakan keluarga ChaCha (ChaCha lainnya: 

Chacha6, ChaCha7) 

• Alasannya terutama adalah untuk memperoleh difusi yang lebih baik pada 

struktur internal sekaligus mempertahankan kecepatan dan kesederhanaan 

desain ARX (add, rotate, xor).

• ChaCha20 Mempertahankan keamanan Salsa20 sambil memperbaiki karakteristik 

desain

• Salah satu tujuan penting desain ChaCha adalah memperoleh implementasi yang 

sangat efisien pada perangkat lunak.

• ChaCha20 menjadi sangat populer dalam protokol modern

90

\end{document}
